\begin{prob}
    Write each of the following functions in $\Theta$-notation in simplest terms.
    You do not need to show your work or provide justification.

    Hint: your answer should have the form $\Theta(n^\alpha)$, $\Theta(n^\alpha
    \gamma^{\delta n})$, $\Theta((\log_\mu(n))^\beta)$, $\Theta(\gamma^{\delta
    n})$, or $\Theta(1)$, where $\alpha, \beta, \gamma, \delta, \mu$ are
    constants that depend on the problem. While you \emph{can} find the answer
    using more formal means (such as limits), you should be able to find the
    answers rather quickly using informal reasoning. You can always use the
    formal approach to check your answer, if you wish.

    \begin{subprobset}

        \begin{subprob}
            $\displaystyle f(n) = \sqrt{\frac{n^6 + n^5 \log n - 1000n}{n^3 - n^2}}$

            \begin{soln}
                $f(n) = \Theta(n^{3/2})$
            \end{soln}
        \end{subprob}

        \begin{subprob}
            $f(n) = (n + 1)(n-1) \times 2^{\log_2 (3^n) }$

            \begin{soln}
                $f(n) = \Theta(n^2 \cdot 3^n)$
            \end{soln}
        \end{subprob}

        \begin{subprob}
            $f(n) = 2^{1000} + 10 + \log_{10} 1000$

            \begin{soln}
                $f(n) = \Theta(1)$.
            \end{soln}
        \end{subprob}

        \begin{subprob}
            $f(n) = \log_2(3^{(n^2 + 2n)} - 5n^3 + \log_2 n)$

            \begin{soln}
                $f(n) = \Theta(n^2)$.

                $3^{(n^2 + 2n)}$ grows a lot faster than the other terms inside of the logarithm,
                so we can ignore them. Therefore we ask: what is $\log_2(3^{(n^2 + 2n)})$ simplified
                in $\Theta$-notation?

                You might remember that when you take the log of an exponential, the exponent
                comes down in front of the log. To be precise,
                if you look up a table of logarithm properties, you'll find
                the identity $\log_a b^x = x \log_a b$.  Therefore $\log_2(3^{(n^2 + 2n)}) = (n^2 + 2n) \log_2 3 = \Theta(n^2)$.

                In general, $\log_a b^{g(n)} = \Theta(g(n))$, as long as $a$ and $b$ are constants independent of $n$
                and greater than one.
            \end{soln}
        \end{subprob}

        \begin{subprob}
            $f(n) = 2^n + 2^{10n} $

            \begin{soln}
                $f(n) = \Theta(2^{10n})$

                Informally,
                $2^{10n}$ grows a lot faster than $2^{n}$, so this is dominated by $2^{10n}$.
                Although we often ignore constants when working with asymptotic notation,
                we have to be careful -- we can't \emph{always} ignore them!
                Exponents are one place where constants \emph{do} matter. If you're ever
                unsure, go back to the definition of $\Theta$-notation and try to prove
                that the function is $\Theta$ of what your intuition tells you, as this
                will let you know whether you can indeed ignore constants.

                If you want to see this more formally, the easiest way might be to use
                limits. We'll show that $2^{10n}$ is much larger than $2^n$ asymptotically
                by taking the limit
                \[
                    \lim_{n \to \infty} \frac{2^{10n}}{2^n}
                \]
                and showing that it is $\infty$.

                \begin{align*}
                    \lim_{n \to \infty}
                        \frac{2^{10n}}{2^n}
                    &=
                    \lim_{n \to \infty}
                        2^{(10 - 1)n}
                        \\
                    &=
                    \lim_{n \to \infty}
                        2^{9n}
                        \\
                    &= \infty
                \end{align*}

                Therefore, $2^{10n}$ is \emph{not} $\Theta(2^n)$.
            \end{soln}
        \end{subprob}


        \begin{subprob}
            $f(n) = \log_2(n^2 + 3n + 5) \times \log_2(3n^3)$

            \begin{soln}
                $\Theta((\log_2 n)^2)$. The base doesn't actually matter, though, so
                we typically write $\Theta((\log n)^2)$.
            \end{soln}
        \end{subprob}

    \end{subprobset}
\end{prob}
